\section{Weak unification: four routes and four lines}\label{sec:debil}
\subsection{Route 1: the monoidal defect as integration}\label{sec:via1}
Let $X=(A,E)$ be a system of $\mathbf{Open}_G$ with stationary covariance $P$. For an $n\times n$ matrix $B$, used as a linear predictor of $X_{t+1}$ from $X_t$ with residual covariance $E+(A-B)P(A-B)^\top$, let $I_B=\tfrac12\log\det P-\tfrac12\log\det\bigl(E+(A-B)P(A-B)^\top\bigr)$ be the predictive information that $B$ captures. With $B=A$ one has $I_A=I_{\rm pred}=I(X_t;X_{t+1})$.
\begin{proposition}[Extensive decomposition plus defect]\label{prop:decomposition}
For systems with at least two units, $I_{\rm pred}=I_{\rm cut}+J$, with $I_{\rm cut}=\max_\pi\max_{B\in\mathfrak B_\pi}I_B$. (With one unit there is no bipartition, $J=+\infty$ by the convention of Remark~\ref{rem:small} and $I_{\rm cut}=\max\varnothing=-\infty$, so the identity is undefined; on the empty system neither $J$ nor $I_{\rm cut}$ is defined.) Moreover:
\begin{itemize}
\item $I_{\rm pred}$ is additive under $\otimes$, and it vanishes on the empty system;
\item $J$ vanishes exactly on the systems with at least two units and disconnected support (Lemma~\ref{lem:J}), which form a tensor ideal $\mathcal Z$: $X\in\mathcal Z$ implies $X\otimes Y\in\mathcal Z$ for every $Y$. $\mathcal Z$ contains the decomposable objects $X\otimes Y$ ($X,Y$ nonempty) and coincides with them only when $E$ is block diagonal along a cut without crossing coefficients; with noise correlated between the blocks, $J=0$ without the system being a tensor product (Section~\ref{sec:l2});
\item for nonempty $X,Y$, $I_{\rm cut}(X\otimes Y)=I_{\rm pred}(X)+I_{\rm pred}(Y)$; if $X$ and $Y$ have at least two units, this equals $I_{\rm cut}(X)+I_{\rm cut}(Y)+J(X)+J(Y)\ge I_{\rm cut}(X)+I_{\rm cut}(Y)$, with defect $J(X)+J(Y)$.
\end{itemize}
\end{proposition}
\begin{proof}
By the definition of $J$, the minimum of $\tfrac12\log\det\bigl(E+(A-B)P(A-B)^\top\bigr)$ over the bipartitions $\pi$ and the predictors $B\in\mathfrak B_\pi$ is $\tfrac12\log\det E+J$, so $I_{\rm cut}=\tfrac12\log\det P-\tfrac12\log\det E-J=I_{\rm pred}-J$. The additivity of $I_{\rm pred}$ is Proposition~\ref{prop:monoidal}; on the empty system both determinants are empty and equal to $1$. The vanishing set of $J$ is given by Lemma~\ref{lem:J}. If a bipartition $(V_1,V_2)$ of the units of $X$ is crossed by no coefficient, then so is the bipartition $(V_1\cup V_Y,V_2)$ of the units of $X\otimes Y$, where $V_Y$ are the units of $Y$; and for nonempty $X,Y$ the bipartition into the units of $X$ and those of $Y$ is crossed by none. Finally, for nonempty $X,Y$, $J(X\otimes Y)=0$ (Proposition~\ref{prop:monoidal}), so $I_{\rm cut}(X\otimes Y)=I_{\rm pred}(X\otimes Y)=I_{\rm pred}(X)+I_{\rm pred}(Y)$, and $I_{\rm pred}=I_{\rm cut}+J$ for $X$ and $Y$ with at least two units gives the rest.
\end{proof}
$J$ is the minimum divergence, over the cuts, from the manifold of disconnected predictors, which is the information-geometric reading~\cite{Oizumi2016PNAS}. The projections of IIT and GNWT type factor through the operation that sends $\mathcal Z$ to zero. The expression ``exact sequence'' can be used here only as an additive splitting in the monoid $(\mathbb R_{\ge0},+)$: there is no abelian structure. Calling $I_{\rm pred}$ the ``FEP base term'' is an interpretation: what is proved is extensivity.

\begin{theorem}[Min-cut sandwich]\label{thm:cheeger}
Let $X=(A,E)$ have at least two units, let $w_{\min}=\min_\pi\lVert A_{\rm cross}(\pi)\rVert_F^2$, where $A_{\rm cross}(\pi)$ keeps only the coefficients of $A$ that cross $\pi$, and let $\kappa_E=\lambda_{\max}(E)/\lambda_{\min}(E)$. Then
\begin{equation}
\tfrac12\log\bigl(1+w_{\min}/\kappa_E\bigr)\ \le\ J\ \le\ \tfrac12\,\frac{\lambda_{\max}(P)}{\lambda_{\min}(E)}\,w_{\min}.
\end{equation}
\end{theorem}
\begin{proof}
\emph{Lower bound.} For every bipartition $\pi$ and every $B\in\mathfrak B_\pi$, with $M=(A-B)P(A-B)^\top$, the proof of Lemma~\ref{lem:J} gives $\det(E+M)/\det E\ge1+\lambda_{\min}(P)\lambda_{\max}(E)^{-1}\lVert A-B\rVert_F^2$. Moreover, $\lambda_{\min}(P)\ge\lambda_{\min}(E)$, and $\lVert A-B\rVert_F\ge\lVert A_{\rm cross}(\pi)\rVert_F$ because $B$ vanishes at the crossing positions. Hence the ratio is at least $1+\lVert A_{\rm cross}(\pi)\rVert_F^2/\kappa_E\ge1+w_{\min}/\kappa_E$.
\emph{Upper bound.} Take a cut $\pi^*$ that attains $w_{\min}$ and $B=A-A_{\rm cross}(\pi^*)\in\mathfrak B_{\pi^*}$, so that $M=A_{\rm cross}(\pi^*)PA_{\rm cross}(\pi^*)^\top$. Since $\log\det(I+Y)\le\operatorname{tr}Y$ for $Y\succeq0$, one gets $J\le\tfrac12\operatorname{tr}(E^{-1}M)\le\tfrac12\lambda_{\min}(E)^{-1}\lambda_{\max}(P)\lVert A_{\rm cross}(\pi^*)\rVert_F^2$.
\end{proof}
It was checked numerically on the four systems and on twelve random systems with random $E$, of which eleven have connected support; the remaining one, the first drawn, has no off-diagonal coupling, so both bounds and $J$ vanish (\texttt{pistas\_numerico.py}). It resembles the Cheeger inequality in that it bounds a global quantity by a cut, but no spectral gap or algebraic connectivity is involved. Efficient searches for the minimum-information bipartition of geometric integrated information exist: Queyranne's algorithm for symmetric submodular functions, applied although the measure is not submodular~\cite{Kitazono2018}, and spectral clustering of the correlation matrix~\cite{TokerSommer2019}.

\emph{Status.} \stA\ \stM. Besides the paper proof above, the theorem is checked in Lean with Mathlib exactly as stated, with $\lambda_{\min}$ and $\lambda_{\max}$ the extreme eigenvalues, for $n\ge2$, $E\succ0$ and $P\succ0$ given as data with $P=APA^\top+E$ (\texttt{QAudit.sandwich}). A version with arbitrary Loewner bounds in place of the eigenvalues, $e_1I\preceq E\preceq e_2I$ with $e_1>0$ and $P\preceq p_2I$, is also checked (\texttt{QAudit.sandwich\_lower}, \texttt{QAudit.sandwich\_upper}); the upper bound needs only $P\succeq0$, not the stationarity equation. With the stationary covariance constructed as the series $\sum_{j\ge0}A^jE(A^j)^\top$ when $\operatorname{spr}(A)<1$, the sandwich is \texttt{QAuditBridge.sandwich\_spr}. As in Lemma~\ref{lem:J}, the inner minimum is an infimum whose attainment is checked separately, and the case of one unit, with $w_{\min}=+\infty$ (Remark~\ref{rem:small}), is not formalized.

\subsection{Route 2: a lax 2-cell between broadcast and integration}\label{sec:via2}
The target is enriched over the poset $\mathbf 2$: on $\mathbf{Open}_G^{+}$, a 2-cell $\alpha:G_{\rm GNWT}F_{\rm GNWT}\Rightarrow G_{\rm IIT}F_{\rm IIT}$ is the pointwise inequality.
\begin{theorem}[Lax 2-cell]\label{thm:laxa}
\begin{itemize}
\item[(a)] A broadcast hub implies connectivity of the undirected graph of the support (\texttt{universal\_connected}, for a symmetric adjacency) for every $n\ge1$; by Lemma~\ref{lem:J} for $n\ge2$, and by Remark~\ref{rem:small} for $n=1$, it implies $J>0$.
\item[(b)] For $1\le n\le3$, broadcast and connectivity coincide: \texttt{n3\_equiv} for $n=3$, \texttt{n2\_equiv} for $n=2$, and $n=1$ is trivial.
\item[(c)] For every $n\ge4$, the path $P_n$ is connected (\texttt{path\_reach\_from\_zero}) and has no hub (\texttt{path\_no\_universal}): the 2-cell is strict for \emph{every} $n\ge4$, not only for $n\in\{4,5\}$. For $n=4$ there are exactly $15$ connected labeled graphs without a hub (\texttt{n4\_count}): the $12$ labelings of $P_4$ and the $3$ of $C_4$, exactly the two types of the numerical example, while the star $K_{1,3}$ has a hub. For $n=5$ there are $472$: they are the $728$ connected labeled graphs minus the $256$ with a hub, by inclusion--exclusion. This count is not in Lean; it is checked by exhaustive enumeration of the $2^{10}$ labeled graphs on five units in \texttt{pistas\_numerico.py}, which also checks the counts $0$, $0$ and $15$ for $n=2,3,4$.
\item[(d)] Quantitative version, for $n\ge2$: if a unit $w$ satisfies $\max(|A_{vw}|,|A_{wv}|)\ge c$ for every $v\ne w$ (a hub of strength $c$), then $J\ge\tfrac12\log(1+c^2/\kappa_E)$. If the coupling has strength $c$ in both directions, $|A_{vw}|,|A_{wv}|\ge c$, the bound is $J\ge\tfrac12\log(1+2c^2/\kappa_E)$.
\end{itemize}
\end{theorem}
\begin{proof}[Proof of (d)]
Every cut separates the hub $w$ from some unit $v$, so $A_{vw}$ and $A_{wv}$ cross it. Hence $w_{\min}\ge c^2$, or $w_{\min}\ge2c^2$ if both have strength $c$. Theorem~\ref{thm:cheeger} applies.
\end{proof}
\emph{Scope.} The 2-cell holds when both projections are computed on \emph{the same} causal model and at the same grain. It is not a theorem about the theories as such. IIT and GNWT can also disagree when the grain or the substrate changes: for example, a workspace architecture implemented on a conventional computer would satisfy functional broadcast while having, according to IIT, a negligible $\Phi$~\cite{TononiKoch2015}. Moreover, $J$ is not the $\Phi$ of IIT~3.0.

\subsection{Route 3: selection in the fibres}\label{sec:via3}
\begin{theorem}[\texttt{section\_factor\_iff}]\label{thm:section}
Let $\Phi:D\to V$ and $u:D\to U$ be functions, and let $s:V\to D$ be a section of $\Phi$ on its image, that is, $\Phi(s(\Phi(x)))=\Phi(x)$ for every $x\in D$. Then the values at the representatives define a factor of $u$, that is, $u=u\circ s\circ\Phi$, if and only if $S_u$ already holds: $\Phi(x)=\Phi(y)\Rightarrow u(x)=u(y)$.
\end{theorem}
With $\Phi=\Phi_Q$ on $D=D_0$, choosing representatives adds nothing to $S_u$: the other realizations in the fibre remain in $D_0$ and must have the same degree. The version with content is a \emph{physical selection law}.
\begin{definition}[SEL]\label{def:SEL}
Let $\mathcal F$ be a real-valued functional on $D_0$. The \emph{physical selection law} SEL for $\mathcal F$ states that every realization in $D_0$ belongs to the unique class, up to isomorphism, that minimizes $\mathcal F$ in its fibre: for every $s\in D_0$, the minimizers of $\mathcal F$ on the fibre $\{t\in D_0:\Phi_Q(t)=\Phi_Q(s)\}$ form a single isomorphism class at grain $\gamma$, and $s$ belongs to it.
\end{definition}
Under SEL, any two realizations in $D_0$ with the same value of $\Phi_Q$ belong to the minimizing class of their common fibre, so every fibre of $\Phi_Q$ on $D_0$ is a single isomorphism class (fibre purity). With OI$_\gamma$ (Definition~\ref{hyp:OI}), which makes $u$ constant on isomorphism classes, purity gives $S_u$ (\texttt{purity\_gives\_Su}). But SEL contradicts Rob as soon as $\Phi_Q$ is continuous there. Assume Rob (Definition~\ref{hyp:Rob}), with the conscious realization $s_0$, the convex set $P$ of kernel parameters and the realizations $s_\theta$ as there, and suppose that $\theta\mapsto\Phi_Q(s_\theta)$ is continuous on $P$ and that $\Phi_Q(s_0)$ is interior to its range. Then Proposition~\ref{prop:extension}(i), which rests on Theorem~\ref{thm:continuum}, gives uncountably many pairwise non-isomorphic realizations in the fibre through $s_0$, and a single non-isomorphic pair in one fibre already destroys purity (\texttt{witness\_breaks\_purity}). Without continuity there is no contradiction: an admissible descriptor that is discontinuous and injective on the classes makes all fibres pure.

Three points about FEP in this role:
\begin{itemize}
\item variational free energy is defined for each system relative to its own generative model, so comparing different systems requires a common reference that is not given;
\item the non-equilibrium steady state (NESS) is a state of a system, not a selection among systems;
\item the uniqueness of the minimizer is not guaranteed.
\end{itemize}

\subsection{Route 4: the connected domain}\label{sec:via4}
On $\mathcal C_{J>0}=\{C_4,K_{1,3},P_4\}$ the integration presence predicate equals $1$ everywhere, so every presence 2-cell into it is trivial (\texttt{connected\_domain\_trivial}). The non-commutation with broadcast persists on $C_4$ and $P_4$. Restricting the domain does not remove it: it removes the case in which both theories say 0. Only the quantitative 2-cell (Theorem~\ref{thm:laxa}(d)) provides information on $\mathcal C_{J>0}$.

\subsection{Weak unification theorem (structural)}\label{sec:unif}
\begin{theorem}[Structural unification in $\mathbf{Open}_G$]\label{thm:unif}
In the symmetric monoidal groupoid of open linear-Gaussian systems, with structural broadcast and blankets:
\begin{enumerate}
\item[(i)] \textbf{Extensive plus defect.} For systems with at least two units, $I_{\rm pred}=I_{\rm cut}+J$, and $J>0$ if and only if the support is connected. $I_{\rm pred}$ is strongly monoidal, and $J$ and broadcast vanish on the tensor ideal $\mathcal Z$.
\item[(ii)] \textbf{Sandwich.} For systems with at least two units, $J$ is comparable to $w_{\min}$ in the sense of Theorem~\ref{thm:cheeger}.
\item[(iii)] \textbf{Lax 2-cell.} A hub of strength $c$ gives $J\ge\tfrac12\log(1+c^2/\kappa_E)$, and $J\ge\tfrac12\log(1+2c^2/\kappa_E)$ if the coupling has that strength in both directions. The 2-cell is strict for every $n\ge4$ and is an equality of predicates for $1\le n\le3$.
\item[(iv)] \textbf{Retained obstructions.}
  \begin{itemize}
  \item There is no strict commutation of the presence bridges.
  \item There is no translation between positive integration and broadcast.
  \item There are blind spots.
  \item No selection in the fibres gives $S_u$ without changing the domain, and the law SEL (Definition~\ref{def:SEL}) contradicts Rob if $\theta\mapsto\Phi_Q(s_\theta)$ is continuous on the convex set $P$ of Rob and $\Phi_Q(s_0)$ is interior to its range (Section~\ref{sec:via3}). This point concerns the original domain $D_0$, not $\mathbf{Open}_G$.
  \end{itemize}
\end{enumerate}
\end{theorem}
\emph{Status.} (i)--(iii) are theorems about invariants of the model, on the physical side, with numerical checks. Their analytic statements about $J$, that $J>0$ if and only if the support is connected, so that $J$ vanishes exactly on $\mathcal Z$, and the sandwich, are proved on paper and checked in Lean with Mathlib (Lemma~\ref{lem:J}, Theorem~\ref{thm:cheeger}) \stA\ \stM. The decomposition $I_{\rm pred}=I_{\rm cut}+J$, the strong monoidality of $I_{\rm pred}$ and the hub bound in (iii) are proved on paper only \stA, and the combinatorial parts are checked in Lean without Mathlib \stL. (iv) collects the compiled counterexamples. They order the projections into a graded family with a common defect structure. They relate two structural proxies computed on the same model, not IIT and GNWT as theories (Section~\ref{sec:via2}, \emph{Scope}). They establish no bridge law, and $S_u$ remains open.


{\sloppy Lean without Mathlib checks only the logical and combinatorial skeleton of this block, on Boolean predicates and abstract types \stL. \texttt{FiltracionCociente.lean} contains \texttt{noncomm\_locus}, \texttt{commute\_off\_band}, \texttt{zmul\_assoc}, \texttt{strong\_monoidal\_zero}, \texttt{factors\_through\_collapse}, \texttt{restriction\_vacuous}, \texttt{seg\_is\_Su}, \texttt{seg\_implies\_Su}, \texttt{segd\_iff} and \texttt{pullback\_determined}. \texttt{PistasUnificacion.lean} contains \texttt{universal\_connected}, \texttt{n3\_equiv}, \texttt{path\_reach\_from\_zero}, \texttt{path\_no\_universal}, \texttt{n4\_count}, \texttt{tensor\_no\_universal}, \texttt{section\_factor\_iff}, \texttt{purity\_gives\_Su}, \texttt{witness\_breaks\_purity} and \texttt{connected\_domain\_trivial}; \texttt{n2\_equiv} is in \texttt{GaussConexion.lean}. \texttt{noncomm\_locus} and \texttt{commute\_off\_band} take as a hypothesis that the hub implies connectivity. \texttt{strong\_monoidal\_zero} assumes strict unit laws for the tensor and that every object other than the unit is nonempty, and, on the projection, the three conditions stated in Section~\ref{sec:l2}: the unit goes to the adjoined unit $e$, every nonempty object receives a value other than $e$, and every product of two nonempty objects has the value $0$. Lean with Mathlib checks the following analytic statements about $J$, for $n\ge2$ \stM: Lemma~\ref{lem:J} and the vanishing of $J$ exactly on disconnected supports (\texttt{QAudit.lemmaJ\_stationary}, \texttt{QAudit.J\_eq\_zero\_iff}); the bounds on $J$ (Theorem~\ref{thm:cheeger}, \texttt{QAudit.sandwich}); and the identification of the presence $J>0$ with $w_{\min}>0$ (\texttt{QAuditBridge.J\_pos\_iff\_wmin\_pos}), which, with $J\ge0$, gives the vanishing of $J$ on the class $\mathcal I$ defined in Section~\ref{sec:l2}. The following remain on paper \stA: the decomposition $I_{\rm pred}=I_{\rm cut}+J$, the additivity of $I_{\rm pred}$, the hub bound of Theorem~\ref{thm:laxa}(d) and what is said about the restricted Gaussian family. The conditional unification proposition is not formalized.\par}

\subsection{Line 1: filtration by connectivity}\label{sec:l1}
Let $\mathcal C_{\rm Central}$ be the systems with a broadcast hub and $\mathcal C_{w_{\min}>0}$ the nonempty systems with connected support, that is, with $w_{\min}>0$: the one-unit systems, for which $w_{\min}=+\infty$ (Remark~\ref{rem:small}), and the systems with at least two units and $w_{\min}>0$. On nonempty systems, $w_{\min}>0$ is equivalent to $J>0$, by Lemma~\ref{lem:J} and Remark~\ref{rem:small}; for $n\ge2$ this equivalence is checked in Lean with Mathlib (\texttt{QAuditBridge.J\_pos\_iff\_wmin\_pos}) \stM. There is a filtration by full subcategories
\begin{equation}
\mathcal C_{\rm Central}\ \subseteq\ \mathcal C_{w_{\min}>0}\ \subseteq\ \mathbf{Open}_G^{+},
\end{equation}
whose first inclusion is \texttt{universal\_connected} (for one unit it is trivial). The complement of $\mathcal C_{w_{\min}>0}$ in $\mathbf{Open}_G^{+}$ is the class $\mathcal I$ of Section~\ref{sec:l2}, the systems with at least two units and disconnected support; the empty system lies outside $\mathbf{Open}_G^{+}$ (Remark~\ref{rem:small}). For every $n\ge4$ the first inclusion is strict: the path $P_n$ is connected (\texttt{path\_reach\_from\_zero}) and has no hub (\texttt{path\_no\_universal}).
\begin{theorem}[Exact locus of non-commutation; \texttt{noncomm\_locus}, \texttt{commute\_off\_band}]\label{thm:locus}
On $\mathbf{Open}_G^{+}$, the presence bridges of IIT type ($J>0$) and GNWT type (hub) disagree exactly on the band $\mathcal C_{w_{\min}>0}\setminus\mathcal C_{\rm Central}$, and commute outside it.
\end{theorem}
\begin{proof}
Apply \texttt{noncomm\_locus} and \texttt{commute\_off\_band} to the Boolean predicates ``has a hub'' and ``$J>0$''; on nonempty systems the latter is connectivity of the support (Lemma~\ref{lem:J} and Remark~\ref{rem:small}). Their hypothesis, that a hub implies connectivity, is Theorem~\ref{thm:laxa}(a).
\end{proof}
For $n=4$ the band has $15$ labeled graphs: the paths $P_4$ and the cycles $C_4$. The lax 2-cell is the inclusion of the filtration. At the central level it is quantitative: a hub of strength $c$ gives $J\ge\tfrac12\log(1+c^2/\kappa_E)$, or $\tfrac12\log(1+2c^2/\kappa_E)$ with coupling in both directions (Theorem~\ref{thm:laxa}(d)). On the band, whose systems have at least four units, only the minimum-cut bound holds, $J\ge\tfrac12\log(1+w_{\min}/\kappa_E)>0$. No global equivalence is possible: $P_n$ is in the band for every $n\ge4$.

\subsection{Line 2: quotient by the factorizable ideal}\label{sec:l2}
$\mathcal I=\{X:\ X\text{ has at least two units and }w_{\min}(X)=0\}$ is the class of systems with at least two units and disconnected support; by Lemma~\ref{lem:J} it is the ideal $\mathcal Z$ of Section~\ref{sec:via1}. A one-unit system is not in $\mathcal I$, since $w_{\min}=+\infty$ there, and neither is the empty system, on which $w_{\min}$ is not defined (Remark~\ref{rem:small}). $\mathcal I$ contains the decomposable objects $X\otimes Y$ ($X,Y$ nonempty), but it does not coincide with them if $E$ is not block diagonal along the cut: with $A=\mathrm{diag}(0.5,\,0.4)$ and $E=\bigl(\begin{smallmatrix}1&0.6\\0.6&1\end{smallmatrix}\bigr)$ one has $w_{\min}=J=0$, and the stationary mutual information between the two units is $0.2188$ nats, so the system is not a tensor product. ``Factorizable'' here means factorizable in the dynamics $A$, not in the joint law. It is a tensor ideal: $X\in\mathcal I\Rightarrow X\otimes Y\in\mathcal I$, and $X\otimes Y\in\mathcal I$ for nonempty $X,Y$. The collapse $q:\mathbf{Open}_G\to\mathbf{Open}_G/\mathcal I$ sends $\mathcal I$ to a point. Since $\mathbf{Open}_G$ is a groupoid, the collapse is not an abelian-category quotient.
\begin{theorem}[\texttt{factors\_through\_collapse}, \texttt{strong\_monoidal\_zero}]\label{thm:collapse}
\begin{itemize}
\item[(a)] A projection factors through $q$ if and only if it is constant on $\mathcal I$. The projections of IIT type (the value of $J$ and the presence $J>0$) and broadcast are constant on $\mathcal I$, since they take the value $0$ on $\mathcal I$. $I_{\rm pred}$ is not.
\item[(b)] Precise replacement of monoidality: let $(D\sqcup\{e\},\circledast)$ have unit $e$ and $a\circledast b=0$ for $a,b\ne e$, where $0\in D$ is a fixed zero value; it is a monoid (\texttt{zmul\_assoc}). If one adopts the convention that the empty system goes to the unit $e$ and no nonempty system goes to $e$, then $J$, with $D=[0,+\infty]$, and broadcast, with $D=\{0,1\}$ for the presence of a hub or $D=\mathbb N$ for the number of hubs (zero $0$), are \emph{strong monoidal} functors into that zero-product monoid: $F(\varnothing)=e$ and $F(X\otimes Y)=F(X)\circledast F(Y)$ for all $X,Y$. Without that convention they are not: if the empty system went to a value of $D$, such as $0$, then for $X=K_{1,3}$ one would have $F(X\otimes\varnothing)=F(X)\ne0=F(X)\circledast F(\varnothing)$. This is a reformulation of the fact that they vanish on the products of nonempty systems, not an additional property. The set of hubs itself is not an isomorphism invariant, so by Proposition~\ref{prop:funtor} it is not a functor into this discrete monoid; it is transported along isomorphisms, as a functor into finite sets and bijections (Section~\ref{sec:tesis}). $I_{\rm pred}$ is strong monoidal into $(\mathbb R_{\ge0},+,0)$.
\end{itemize}
\end{theorem}
\begin{proof}
(a) is \texttt{factors\_through\_collapse}, with the point to which $\mathcal I$ collapses sent to the constant value; the values on $\mathcal I$ are given by Lemma~\ref{lem:J} for $J$ and by Theorem~\ref{thm:laxa}(a) for broadcast, since a system with a hub has connected support. Some systems of $\mathcal I$ have different values of $I_{\rm pred}$, for instance $A=\mathrm{diag}(a,a)$ and $E=I$ with different values of $a$. (b) Lean proves the abstract form, \texttt{strong\_monoidal\_zero}, for a tensor, a unit and a predicate ``nonempty'' on an abstract type of objects, under six hypotheses: strict left and right unit laws, $\mathrm{unit}\otimes x=x$ and $x\otimes\mathrm{unit}=x$ (\texttt{hunitL}, \texttt{hunitR}); every object other than the unit is nonempty (\texttt{hne}); the projection sends the unit to $e$ (\texttt{hF0}); every nonempty object receives a value in $D$, not $e$ (\texttt{hFne}); and every product of two nonempty objects has the value $0$ (\texttt{hvan}). The first three concern the monoidal structure: they hold on isomorphism classes of systems, on which an isomorphism-invariant projection is defined and on which the class of the empty system is the only one without units. \texttt{hF0} and \texttt{hFne} are the convention; for $J$, \texttt{hFne} holds because $J=+\infty$ on one-unit systems (Remark~\ref{rem:small}). \texttt{hvan} is Proposition~\ref{prop:monoidal}: $J(X\otimes Y)=0$, and $X\otimes Y$ has no hub (\texttt{tensor\_no\_universal}). The additivity of $I_{\rm pred}$ is Proposition~\ref{prop:monoidal}.
\end{proof}
\emph{Volume and obstruction:} for systems with at least two units, $I_{\rm pred}=I_{\rm cut}+J$, where $I_{\rm pred}$ is the extensive term, which survives in $X\otimes Y$, and $J$ is the obstruction to partition. The correct bound is $J\ge\tfrac12\log(1+w_{\min}/\kappa_E)$, not $J\ge w_{\min}$. Calling $I_{\rm pred}$ the ``FEP term'' remains an interpretation (Section~\ref{sec:via1}).

\subsection{Line 3: variational selection in the fibres}\label{sec:l3}
\begin{theorem}[\texttt{section\_factor\_iff}, \texttt{restriction\_vacuous}, \texttt{segd\_iff}]\label{thm:selection}
\begin{itemize}
\item[(a)] No section of $\Phi$, not even the section that chooses a free-energy minimizer, gives $S_u$: the values at the representatives define a factor if and only if $S_u$ already holds.
\item[(b)] Restricting the domain to the image of a section leaves each fibre with a single element, and $S_u$ becomes trivially true for \emph{every} $u$. The obligation is not met: it is eliminated by changing the domain, which this audit excludes by method: it does not change the descriptor, the domain, the components or the meaning of Q (Section~\ref{sec:intro}).
\item[(c)] Requiring the fibre to ``satisfy SEG'' in order to collapse it is circular. Under $B_\psi$, SEG (Definition~\ref{def:SEG}) holds if and only if $S_u$ holds with $u=\psi\circ\mathfrak U$ nonconstant on $D_0$ (\texttt{segd\_iff}: for the right-to-left direction, take the level sets of $\psi$ as the groups), and already a split into at least two groups with $\psi$ constant on each group and condition~(ii) of Definition~\ref{def:SEG} implies $S_u$ (\texttt{seg\_implies\_Su}). \texttt{seg\_is\_Su} only records that ``equal $\Phi$, equal $\psi\circ\mathfrak U$'' is, by definition, $S_u$ (\texttt{Iff.rfl}).
\end{itemize}
\end{theorem}
\emph{The restricted Gaussian family} of Section~\ref{sec:gauss} is an auxiliary test bench. Its descriptor is $J$, which is not $\Phi_Q$, and its obligation is $S_J$, which is not $S_u$. In this family, at the level of laws, each fibre of $J$ is a single class: the normalization $Z_i=X_i/\sqrt{s_i}$ gives identity noise and coupling $\sqrt k$, with $k=c^2s_2/s_1$, so there is nothing to select. The multiplicity lies in the pointed state, $z\sim N(0,\mathrm{diag}(1+k,1))$, and in the histories $z_{0:T}$. If the free energy of a state is taken to be its surprisal $-\log p_\theta(z)$, choosing the state of minimum free energy, the mode $z=0$, discards the realizations: they have $z\ne0$ almost surely, and $S_J$ requires that they too have the degree of the mode. Finally, a \emph{physical selection} law (only the minimizers are realized) would give fibre purity. In the family it contradicts the model itself, in which $z\ne0$ almost surely; in the original domain it is the law SEL of Definition~\ref{def:SEL}, which contradicts Rob when $\theta\mapsto\Phi_Q(s_\theta)$ is continuous on the convex set $P$ of Rob and $\Phi_Q(s_0)$ is interior to its range (Section~\ref{sec:via3}).

\subsection{Line 4: the weak limit}\label{sec:l4}
The pullback of a section $s$ along $\Phi$ is the subcategory of canonical representatives, $\{X: X=s(\Phi(X))\}$. On it, every additional function of the representative is a function of $\Phi$ (\texttt{pullback\_determined}): the weak limit adds no information to the section. The cut spectrum $w_{\min}$ is one of those functions, but not an invariant of $\mathbf{Open}_G$. Positive rescalings, which are isomorphisms, change it: on $C_4$, $w_{\min}=0.36$ with $D=I$ and $0.765$ with $D=\mathrm{diag}(1,2,1,2)$, with $J=0.2213$ in both cases. Only its vanishing is invariant, and the sandwich of Theorem~\ref{thm:cheeger} holds for each representative. Linking the weak limit to $S_u$ or to SEG therefore presupposes $S_u$: by part~(a) of Theorem~\ref{thm:selection}, the values at the representatives define a factor only if $S_u$ already holds.

The only thing that can be asserted here is conditional and, given $S_u$, immediate:
\begin{proposition}[Conditional unification]\label{prop:conditional}
Assume $B_\psi$. If SEG holds on $D_0$ (that is, $S_u$ with $u_{\rm orig}$ nonconstant; Definition~\ref{def:SEG}), then:
\begin{itemize}
\item $u_{\rm orig}=h\circ\Phi_Q$ with $h$ unique on $\Phi_Q(D_0)$;
\item every section computes $h$;
\item if each realization in $D_0$ is represented by a nonempty system of $\mathbf{Open}_G$ at grain $\gamma$, the filtration by cuts stratifies $D_0$ into the central, band and disconnected-support levels, on which the presence bridges commute (central and disconnected) or disagree (band). This point does not depend on SEG.
\end{itemize}
\end{proposition}
\begin{proof}
SEG gives $S_u$ (Definition~\ref{def:SEG} and \texttt{seg\_implies\_Su}), and $S_u$ is equivalent to the factorization $u_{\rm orig}=h\circ\Phi_Q$ with $h$ unique on $\Phi_Q(D_0)$ (Definition~\ref{def:phiQ}(e)). For a section $s$, \texttt{section\_factor\_iff} gives $u_{\rm orig}=u_{\rm orig}\circ s\circ\Phi_Q$, so $h(v)=u_{\rm orig}(s(v))$ for $v\in\Phi_Q(D_0)$. The third point is Theorem~\ref{thm:locus} applied to the representing systems.
\end{proof}
The premise is the open obligation, not its substitute. Nothing in this block contradicts the compiled counterexamples: no global equivalence is asserted, and $P_n$ always remains in the band.
